\chapter{Единая интенсивная свободная энергия и закрытие детерминантного барьера}

Предыдущая глава показала, что обычный product-KO2 пфаффиан даёт
$-15\log\det R$, тогда как слабый вес $-\frac1{20}\log\det R$ возникает
только после деления логарифма на ранг $300$. Теперь проверим единственную
честную возможность такого деления: один общий интенсивный функционал для
всех классических и квантовых секторов.

\section{Локальная кривизна до нового барьера}

Вдоль бесследовой вариации
\begin{equation}
 R=\frac13I_3+\delta R,
 \qquad \Tr\delta R=0,
\end{equation}
энтропия и классическая цена моста дают
\begin{equation}
 c_{\rm cl}=\frac32+\frac67=\frac{33}{14}.
\end{equation}
Шесть коллективных мостовых флуктуаций дают
\begin{equation}
 c_{\rm br}=-\frac{45}{16},
\end{equation}
поэтому без барьера
\begin{equation}
 c_0=c_{\rm cl}+c_{\rm br}=-\frac{51}{112}<0.
 \label{eq:v6-common-intensive-baseline}
\end{equation}

Обычный пятнадцатикратный пфаффиан добавляет
\begin{equation}
 c_F=\frac92\cdot15=\frac{135}{2}.
 \label{eq:v6-common-intensive-ordinary-fermion}
\end{equation}

\section{Общее деление всей эффективной энергии}

Если определить интенсивную энергию простым общим делением
\begin{equation}
 \mathcal F_{\rm int}=\frac1{300}\mathcal F_{\rm full},
\end{equation}
то знак гессиана не меняется. До деления он равен
\begin{equation}
 c_{\rm full}
 =\frac{33}{14}-\frac{45}{16}+\frac{135}{2}
 =\frac{7509}{112}>0,
\end{equation}
а после деления
\begin{equation}
 \frac{c_{\rm full}}{300}
 =\frac{2503}{11200}>0.
\end{equation}

\begin{proposition}[Общий масштаб не меняет фазу]
Умножение всей эффективной энергии на один положительный коэффициент не
может восстановить отрицательную моду. Поэтому интенсивность как простая
общая единица измерения не спасает детерминантный механизм.
\end{proposition}

\section{Большеразмерный интенсивный предел}

Более содержательная возможность известна из больших-$N$ моделей:
микроскопическое классическое действие является экстенсивным, а свободная
энергия делится на число степеней свободы. Коллективные поля при этом имеют
малое число собственных флуктуаций, и их петли подавлены относительно
экстенсивного седла; см. обзор \path{hep-th/0306133}.

В проекте $B\in M_3(\mathbb R)$ является одним общим семейным мостом. После
удаления орбиты он имеет шесть физических коллективных мод, а не 300
независимых копий. Поэтому единый интенсивный предел даёт
\begin{equation}
 c_{\rm int}
 =\frac{33}{14}
 -\frac1{300}\frac{45}{16}
 +\frac1{300}\frac{135}{2}
 =\frac{5763}{2240}>0.
 \label{eq:v6-common-intensive-large-rank-curvature}
\end{equation}
Изотропное состояние становится строго устойчивым.

Математически это соответствует различию между обычным следом логарифма
и нормированным следовым состоянием. В конечной матричной алгебре
детерминант Фугледе--Кадисона является корнем обычного детерминанта степени
$1/N$; см. \path{arXiv:math/0105109} и \path{arXiv:1107.1059}. Но та же
нормировка должна применяться ко всем операторным логарифмам, включая
мостовой гессиан.

\section{Сколько мостовых копий потребовалось бы}

Пусть существуют $q$ независимых мостовых матриц с тем же отрицательным
гессианом. Тогда
\begin{equation}
 c(q)=\frac{33}{14}
 -\frac{q}{300}\frac{45}{16}
 +\frac9{40}.
\end{equation}
Условие $c(q)<0$ требует
\begin{equation}
 q>\frac{1928}{7}=275\frac37,
\end{equation}
то есть как минимум $276$ независимых копий. Проект содержит одну общую
матрицу $B$, поэтому такая кратность отсутствует. Представительная
кратность физических состояний не создаёт новых интегрируемых
коллективных полей.

\begin{theorem}[Единая интенсивная нормировка закрывает ветвь]
Ни общее деление полной эффективной энергии на $300$, ни стандартный
большеразмерный интенсивный предел не сохраняют одновременно слабый
пфаффиановый барьер и мостовую отрицательную моду. Первый не меняет
отношения коэффициентов, второй подавляет единственную коллективную
мостовую петлю и делает изотропный вакуум устойчивым.
\end{theorem}

\section{Почему выборочная нормировка запрещена}

Если разделить только фермионный логарифм на $300$, но оставить мостовой
детерминант ненормированным, формально получается
\begin{equation}
 \frac{33}{14}-\frac{45}{16}+\frac9{40}
 =-\frac{129}{560}<0.
\end{equation}
Именно этот гибрид дал прежний привлекательный вес $1/20$. Но два
квантовых детерминанта тогда оцениваются разными следами без
операторного основания. Это не общий родительский функционал, а скрытый
относительный коэффициент.

\begin{nogocriterion}[Запрет секторной интенсивности]
Нельзя нормировать фермионный пфаффиан по $M_{300}$, одновременно оставляя
детерминант шести коллективных мостовых мод обычным. Переоткрытие требует
нового доказанного механизма, который физически создаёт разные
экстенсивности секторов, а не выбирает их после проверки знака.
\end{nogocriterion}

\section{Вердикт}

\begin{protocol}
\begin{enumerate}
 \item обычный пфаффиан и прежняя мостовая петля:
 \textbf{изотропный вакуум устойчив};
 \item общее деление всей энергии на $300$:
 \textbf{знак не меняется};
 \item единый большеразмерный интенсивный предел:
 \textbf{$c=5763/2240>0$};
 \item выборочная нормировка только пфаффиана:
 \textbf{даёт отрицательную моду, но запрещена};
 \item минимальное число независимых мостовых копий:
 \textbf{$276$};
 \item число мостовых полей проекта:
 \textbf{одно};
 \item стандартный детерминантный барьер:
 \textbf{закрыт};
 \item мир-кристалл как фазовая картина:
 \textbf{не опровергнут};
 \item следующий допустимый механизм:
 \textbf{негауссова пространственная жёсткость или иной нелокальный
 родительский член};
 \item рождение материи:
 \textbf{не закрыто}.
\end{enumerate}
\end{protocol}

Следующая проверка должна вернуться к уже доказанным пространственным и
суперсвязностным членам Тома V и проверить, может ли производная
жёсткость насыщать проекторный переход до потери ранга без
$-\nu\log\det R$.

Машинный сертификат сохранён в
\path{s2t_v6_common_intensive_free_energy_normalization_gate_results.json}.